% Appendix G: Keyboard Geometry and Motor Fields

\section*{G.1 Introduction}

Conventional computer science treats a keyboard as a discrete symbol emitter.
The gesture framework suggests a different interpretation: \emph{a keyboard
is a geometric field through which hands navigate. Keys are landmarks.
Words are trajectories. Typing is navigation.}

\begin{definition}{Keyboard Geometry}{keyboard-geom}
  A \emph{keyboard geometry} is an embedding $\iota : K \to \mathbb{R}^2$,
  where $K = \{k_1, \ldots, k_n\}$ is the key set.
\end{definition}

\begin{definition}{Finger Assignment and Domains}{finger-assign}
  The \emph{finger assignment map} is $F : K \to \{1,\ldots,8\}$.
  The \emph{finger domain} of finger $i$ is $D_i = F^{-1}(i)$,
  giving a partition $K = D_1 \cup \cdots \cup D_8$.
\end{definition}

\begin{definition}{Stretch Cost}{stretch-cost}
  $\sigma(k) = d(\iota(k),\, h(F(k)))$, where $h(F(k))$ is the home
  position of the responsible finger.
\end{definition}

\begin{definition}{Keyboard Graph}{keyboard-graph}
  $\Gamma = (K,E)$ where $(k_i,k_j) \in E$ whenever a finger may
  move directly between them. Typing is a path $\gamma \subseteq \Gamma$.
\end{definition}

\begin{definition}{Words as Paths}{words-paths}
  A word $w = k_1 k_2 \cdots k_n$ defines a path
  $\gamma_w = (k_1, k_2, \ldots, k_n)$ through $\Gamma$.
\end{definition}

\begin{definition}{Motor Energy of a Word}{motor-energy-word}
  $E(w) = \sum_i d(k_i, k_{i+1})^2$.
  A \emph{geodesic word} minimizes energy among equivalent spellings.
\end{definition}

\begin{definition}{Hand Alternation}{hand-alt}
  With $H(k) \in \{L,R\}$, the \emph{alternation count} is
  $A(w) = \sum_i \mathbf{1}[H(k_i) \neq H(k_{i+1})]$.
  High alternation generally increases typing speed.
\end{definition}

\begin{definition}{Overlap Field}{overlap-field}
  With $a_i(t)$ denoting key activation, $O(t) = \sum_{i<j} a_i(t) a_j(t)$.
  This recovers chord-like information discarded by text systems.
\end{definition}

\begin{definition}{Keyboard Field}{keyboard-field}
  The \emph{keyboard field} is
  $\mathcal{K} = (K, \Gamma, \sigma, H, F)$ — a fully structured
  geometric object, not merely a symbol table.
\end{definition}

\begin{theorem}{Keyboard Navigation Theorem}{keyboard-nav}
  Every typed text corresponds to a path through a constrained motor graph.
\end{theorem}
\begin{proof}
  Each symbol corresponds to a key (vertex of $\Gamma$); successive
  symbols define successive vertices; text determines a path.
\end{proof}

\begin{namedthm}{Symbol Emergence in Keyboard Systems}
  The traditional view: $\text{key} \to \text{symbol}$.
  The geometric view: $\text{trajectory} \to \text{gesture} \to \text{symbol}$.
  Words are not strings. They are routes. Typing is not symbol production.
  It is constrained navigation through a motor field.
\end{namedthm}
